\newpage
\appendix
\section*{Phụ lục A: Mã nguồn Python thực hiện mô phỏng và tối ưu hóa}
\addcontentsline{toc}{section}{Phụ lục A}

Trong phụ lục này, toàn bộ mã nguồn lập trình bằng ngôn ngữ Python trên môi trường Jupyter Notebook được trình bày chi tiết theo từng giai đoạn nghiên cứu của khóa luận.

% =======================================================================
\subsection*{A.1 Tiền xử lý dữ liệu và Thống kê mô tả}
Đoạn mã dưới đây thiết lập các thư viện cấu hình, xử lý chuỗi dữ liệu giá đóng cửa và thực hiện các thống kê mô tả đơn biến, kiểm định phân phối chuẩn (Jarque-Bera).

\begin{lstlisting}[language=Python]
import pandas as pd
import numpy as np
import matplotlib.pyplot as plt
import seaborn as sns
from scipy.optimize import minimize
from scipy.stats import chi2, norm, skew, kurtosis, jarque_bera, gamma
from scipy.stats import t as t_dist

# Cau hinh do thi
sns.set_theme(style="whitegrid")
plt.rcParams['figure.figsize'] = [10, 6]

print("--- Dang xu ly du lieu dau vao ---")
df = pd.read_csv('all_stocks_5yr.csv')
df['date'] = pd.to_datetime(df['date'])

tickers = ['AAPL', 'MSFT', 'AMZN', 'GOOGL']
num_assets = len(tickers)

df_portfolio = df[df['Name'].isin(tickers)].copy()
df_pivot = df_portfolio.pivot(index='date', columns='Name', values='close').dropna()

# Tinh ty suat sinh loi logarit (Log Returns)
returns = np.log(df_pivot / df_pivot.shift(1)).dropna()

# Thong ke mo ta va kiem dinh Jarque-Bera
univariate_stats = pd.DataFrame(index=['Mean (%)', 'Std Dev (%)', 'Skewness', 'Kurtosis', 'JB p-value'])
for t in tickers:
    ret_series = returns[t].values
    jb_stat, jb_p = jarque_bera(ret_series)
    univariate_stats[t] = [
        np.mean(ret_series) * 100,
        np.std(ret_series) * 100,
        skew(ret_series),
        kurtosis(ret_series), 
        jb_p
    ]
print(univariate_stats.round(4))
\end{lstlisting}

% =======================================================================
\subsection*{A.2 Mô hình phụ thuộc và Copula}
Phần này tính toán ma trận tương quan hạng Spearman và hệ số tương quan đuôi dưới thực nghiệm nhằm đánh giá sự phụ thuộc phi tuyến trong trạng thái khủng hoảng.

\begin{lstlisting}[language=Python]
print("\n--- Phan tich Ma tran Tuong quan ---")
print("Ma tran tuong quan Pearson:")
print(returns.corr(method='pearson').round(4))

print("\nMa tran tuong quan hang Spearman:")
spearman_corr = returns.corr(method='spearman')
print(spearman_corr.round(4))

print("\nMa tran he so tuong quan duoi duoi thuc nghiem (Phan vi 5%):")
tail_dep = pd.DataFrame(1.0, index=tickers, columns=tickers)
for i in range(len(tickers)):
    for j in range(i+1, len(tickers)):
        t1, t2 = tickers[i], tickers[j]
        u1 = returns[t1].rank(pct=True)
        u2 = returns[t2].rank(pct=True)
        joint_prob = np.mean((u1 <= 0.05) & (u2 <= 0.05))
        lambda_l = joint_prob / 0.05
        tail_dep.loc[t1, t2] = lambda_l
        tail_dep.loc[t2, t1] = lambda_l
print(tail_dep.round(4))
\end{lstlisting}

% =======================================================================
\subsection*{A.3 Tối ưu hóa danh mục đầu tư Markowitz}
Thuật toán SLSQP được sử dụng để giải bài toán tìm danh mục rủi ro tối thiểu (MVP) và danh mục tiếp điểm có tỷ số Sharpe cực đại, kết hợp vẽ đường biên hiệu quả.

\begin{lstlisting}[language=Python]
print("\n--- Markowitz Portfolio Optimization ---")
num_assets = len(tickers)
mean_returns_ann = returns.mean() * 252
cov_matrix_ann = returns.cov() * 252

# Set risk-free rate to 0 to match original assumption
risk_free_rate = 0.02

def portfolio_performance(weights, mean_returns, cov_matrix):
    p_return = np.sum(mean_returns * weights)
    p_std = np.sqrt(np.dot(weights.T, np.dot(cov_matrix, weights)))
    return p_std, p_return

def negative_sharpe_ratio(weights, mean_returns, cov_matrix, rf):
    p_std, p_return = portfolio_performance(weights, mean_returns, cov_matrix)
    return -(p_return - rf) / p_std

constraints = ({'type': 'eq', 'fun': lambda x: np.sum(x) - 1})
bounds = tuple((0, 1) for asset in range(num_assets))
initial_guess = num_assets * [1. / num_assets,]

sharpe_result = minimize(negative_sharpe_ratio, initial_guess, 
                         args=(mean_returns_ann, cov_matrix_ann, risk_free_rate), 
                         method='SLSQP', bounds=bounds, constraints=constraints)

best_weights = sharpe_result.x
max_sharpe_std, max_sharpe_return = portfolio_performance(best_weights, mean_returns_ann, cov_matrix_ann)
max_sharpe_ratio = max_sharpe_return / max_sharpe_std

print(f"Max Sharpe Ratio: {max_sharpe_ratio:.4f}")
print(f"Expected Annual Return: {max_sharpe_return*100:.2f}%")
print(f"Annual Volatility: {max_sharpe_std*100:.2f}%")
print("Optimal Portfolio Weights:")
for tk, w in zip(tickers, best_weights):
    print(f"  - {tk}: {w*100:.2f}%")

# Plot Efficient Frontier
num_portfolios = 5000
port_results = np.zeros((3, num_portfolios))
np.random.seed(42)
for i in range(num_portfolios):
    w = np.random.random(num_assets)
    w /= np.sum(w)
    p_std, p_return = portfolio_performance(w, mean_returns_ann, cov_matrix_ann)
    port_results[0, i] = p_std
    port_results[1, i] = p_return
    port_results[2, i] = p_return / p_std

plt.figure()
plt.scatter(port_results[0,:], port_results[1,:], c=port_results[2,:], cmap='viridis', marker='o', s=10, alpha=0.3)
plt.colorbar(label='Sharpe Ratio')
plt.scatter(max_sharpe_std, max_sharpe_return, color='red', marker='*', s=200, label='Optimal (Max Sharpe)')
plt.title("Markowitz Efficient Frontier")
plt.xlabel("Risk (Annual Volatility)")
plt.ylabel("Expected Annual Return")
plt.legend()
plt.tight_layout()
plt.savefig('markowitz_efficient_frontier.png')
plt.show()
\end{lstlisting}

% =======================================================================
\subsection*{A.4 Khớp Copula và Tính sai số Cramér-von Mises}
Quá trình ước lượng độ khớp (Goodness-of-Fit) giữa dữ liệu thực nghiệm và các mô hình Copula lý thuyết (Gaussian, t-Student, Clayton) để chọn ra mô hình tối ưu.

\begin{lstlisting}[language=Python]
print("\n--- Copula Goodness-of-Fit and Cramer-von Mises Error ---")
from scipy.stats import gamma

# Tinh toan Du lieu gia quan sat (Pseudo-observations) va Kendall's Tau
pseudo_obs = returns.rank(pct=True)
U = pseudo_obs.values
tau_matrix = returns.corr(method='kendall')
avg_tau = np.mean(tau_matrix.values[np.triu_indices(num_assets, k=1)])
clayton_theta = max(0.1, 2 * avg_tau / (1 - avg_tau))

# Khai bao bac tu do df duoi dang tham so dong
df_val = 4
t_label = f't-Student (df={df_val})'

np.random.seed(42)
sample_idx = np.random.choice(len(U), min(500, len(U)), replace=False)
U_sample = U[sample_idx]

def empirical_copula(u, data):
    return np.mean(np.all(data <= u, axis=1))

def gen_copula_u(c_name, size, df_param):
    mean_0 = np.zeros(num_assets)
    sp_mat = spearman_corr.values
    
    if c_name == 'Gaussian':
        return norm.cdf(np.random.multivariate_normal(mean_0, sp_mat, size))
    elif c_name == t_label:
        w = chi2.rvs(df_param, size=size)
        z = np.random.multivariate_normal(mean_0, sp_mat, size) * np.sqrt(df_param / w)[:, None]
        return t_dist.cdf(z, df=df_param)
    else: 
        v = gamma.rvs(a=1/clayton_theta, scale=1, size=size)
        e = np.random.exponential(scale=1, size=(size, num_assets))
        return (1 + e / v[:, None]) ** (-1/clayton_theta)

gof_scores = {'Gaussian': 0, t_label: 0, 'Clayton': 0}
sim_size_gof = 10000

print("Running calculations, please wait...")
u_sim_dict = {name: gen_copula_u(name, sim_size_gof, df_val) for name in gof_scores.keys()}

for u_pt in U_sample:
    emp_val = empirical_copula(u_pt, U)
    for name in gof_scores.keys():
        gof_scores[name] += (emp_val - empirical_copula(u_pt, u_sim_dict[name]))**2

print("MSE Error Scores (Lower is better):")
for name, score in gof_scores.items():
    print(f"  - {name}: {score:.4f}")

best_copula = min(gof_scores, key=gof_scores.get)
print(f"=> WINNING MODEL: {best_copula.upper()}")
\end{lstlisting}

% =======================================================================
\subsection*{A.5 Mô phỏng Monte Carlo bằng Copula thực nghiệm và Tối ưu hóa CVaR}
Mã nguồn sinh ngẫu nhiên 10.000 kịch bản lợi suất dựa trên mô hình t-Student Copula, sau đó phân tích Giá trị rủi ro (VaR) và thực hiện thuật toán SLSQP để Cực tiểu hóa Giá trị chịu rủi ro có điều kiện (CVaR).

\begin{lstlisting}[language=Python]
print("\n--- Monte Carlo Simulation with t-Student Copula ---")
num_simulations = 10000
df_t = 4 

pearson_approx = 2 * np.sin(np.pi / 6 * spearman_corr.values)
eigenvalues, eigenvectors = np.linalg.eigh(pearson_approx)
eigenvalues[eigenvalues < 0] = 1e-8
pearson_approx_psd = eigenvectors @ np.diag(eigenvalues) @ eigenvectors.T

mean_vector = np.zeros(num_assets)
np.random.seed(42)

Z = np.random.multivariate_normal(mean_vector, pearson_approx_psd, size=num_simulations)
chi2_rvs = np.random.chisquare(df=df_t, size=num_simulations)
W = df_t / chi2_rvs
X = np.sqrt(W)[:, None] * Z
u_sim = t_dist.cdf(X, df=df_t)

sim_returns = np.zeros((num_simulations, num_assets))
for i, ticker in enumerate(tickers):
    sim_returns[:, i] = np.quantile(returns[ticker].values, u_sim[:, i])

# ---- FIGURE 2: EQUAL WEIGHT PORTFOLIO (BLUE) ----
equal_weights = np.array([1.0 / num_assets] * num_assets)
port_sim_returns_eq = np.dot(sim_returns, equal_weights)
var_95_eq = np.percentile(port_sim_returns_eq, 5)

plt.figure(figsize=(10, 6))
sns.histplot(port_sim_returns_eq, bins=50, kde=True, color='skyblue', stat='density', alpha=0.5)
plt.axvline(var_95_eq, color='red', linestyle='--', linewidth=2, label=f'VaR 95% ({var_95_eq*100:.2f}%)')
plt.title(f"Equal Weight Portfolio Return Distribution - t-Student (df={df_t})")
plt.xlabel("Daily Return")
plt.ylabel("Density")
plt.legend()
plt.tight_layout()
plt.savefig('tstudent_sim_returns_var.png')
plt.show()

# ---- CVAR PORTFOLIO OPTIMIZATION (SLSQP) ----
print("\n--- CVaR Portfolio Optimization (SLSQP) ---")
def minimize_cvar(weights, sim_returns, alpha=0.05):
    port_returns = np.dot(sim_returns, weights)
    var_alpha = np.percentile(port_returns, alpha * 100)
    cvar_alpha = port_returns[port_returns <= var_alpha].mean()
    return -cvar_alpha

bounds = tuple((0.0, 1.0) for _ in range(num_assets))
constraints = {'type': 'eq', 'fun': lambda w: np.sum(w) - 1.0}

cvar_opt_result = minimize(
    minimize_cvar, equal_weights, args=(sim_returns, 0.05),
    method='SLSQP', bounds=bounds, constraints=constraints
)

# ---- FIGURE 3: OPTIMAL CVAR PORTFOLIO (PURPLE) ----
if cvar_opt_result.success:
    best_cvar_weights = cvar_opt_result.x
    opt_port_returns = np.dot(sim_returns, best_cvar_weights)
    opt_var_95 = np.percentile(opt_port_returns, 5)
    opt_cvar_95 = opt_port_returns[opt_port_returns <= opt_var_95].mean()
    
    plt.figure(figsize=(10, 6))
    sns.histplot(opt_port_returns, bins=50, kde=True, color='purple', stat='density')
    plt.axvline(opt_var_95, color='red', linestyle='--', linewidth=2, label=f'VaR 95% ({opt_var_95*100:.2f}%)')
    plt.axvline(opt_cvar_95, color='orange', linestyle='-', linewidth=2, label=f'CVaR 95% ({opt_cvar_95*100:.2f}%)')
    plt.title("Optimal Portfolio Return Distribution - Tail Risk Control (Copula-CVaR)")
    plt.xlabel("Daily Return")
    plt.ylabel("Density")
    plt.legend()
    plt.tight_layout()
    plt.savefig('optimal_cvar_portfolio_fixed.png')
    plt.show()
\end{lstlisting}
% =======================================================================
\subsection*{A.6 Mô phỏng Monte Carlo dự báo trong dài hạn (1 năm)}
Quá trình tạo các quỹ đạo tài sản (Asset Paths) trong chu kỳ 252 ngày giao dịch để ước tính hiệu suất và rủi ro ròng vào cuối năm. Lập trình đảm bảo sử dụng hàm t-Student Copula và bộ tỷ trọng tối ưu CVaR.

\begin{lstlisting}[language=Python]
print("\n--- Monte Carlo Simulation Deployment (1 Year / 252 Days) ---")
num_days = 252
portfolio_paths = np.ones((num_days + 1, num_simulations))
np.random.seed(100)

for day in range(num_days):
    Z_day = np.random.multivariate_normal(mean_vector, pearson_approx_psd, size=num_simulations)
    W_day = df_t / np.random.chisquare(df=df_t, size=num_simulations)
    X_day = np.sqrt(W_day)[:, None] * Z_day
    u_day = t_dist.cdf(X_day, df=df_t)
    
    day_returns = np.zeros((num_simulations, num_assets))
    for i, ticker in enumerate(tickers):
        day_returns[:, i] = np.quantile(returns[ticker].values, u_day[:, i])
        
    # Dung bo ty trong toi uu Copula-CVaR de du bao
    portfolio_sim_returns_day = np.dot(day_returns, best_cvar_weights)
    
    # Cong don tai san bang ham mu (Log Returns)
    portfolio_paths[day + 1, :] = portfolio_paths[day, :] * np.exp(portfolio_sim_returns_day)

final_assets = portfolio_paths[-1, :]
var_95_yearly = np.quantile(final_assets, 0.05)
median_yearly = np.median(final_assets)

plt.figure()
plt.plot(portfolio_paths[:, :100], color='blue', alpha=0.08)
plt.axhline(1.0, color='red', linestyle='--', linewidth=1.5)
plt.title("Monte Carlo Simulation Asset Path Scenarios (1 Year)")
plt.savefig('monte_carlo_paths_1year.png')
plt.show()

plt.figure()
sns.histplot(final_assets, bins=50, kde=True, color='seagreen', stat='density')
plt.axvline(var_95_yearly, color='red', linestyle='--', linewidth=2)
plt.title("Distribution of Cumulative Portfolio Asset Value at Year-End")
plt.savefig('final_asset_distribution.png')
plt.show()
\end{lstlisting}